% STATUS: reviewed
\chapter{Diegimas ir eksploatavimas}\label{ch:ops}

Šis skyrius skirtas tam, kas įdiegia \product{} ir palaiko jo veikimą, --- operatoriui. Jis eina ta
tvarka, kuria operatorius susiduria su sistema: kas veikia ir ką turi suteikti serveris, naujas
diegimas, du Compose failo režimai, atnaujinimai, --- o paskui pereina
prie darbų, būdingų būtent šiam produktui: teisės aktų bazės palaikymo, trijų mokamų visos teisės aktų
bazės skaitymų, kuriuos paleidžia tik operatorius, instrukcijų priežiūros ir biudžetų stebėjimo.
Baigiama suplanuotais darbais, būklės patikra ir dažniausiomis nesėkmėmis. Komandos cituojamos
tiksliai, ir jų nedaug; kiekviena vykdoma projekto aplanke serveryje, kuriame neįdiegta nieko,
išskyrus Docker. Skyriuje minimų parametrų ir konstantų sąrašas --- \cref{app:config}.

Tą pačią sritį išsamiau aprašo du saugyklos dokumentai, kurie buvo šio skyriaus žemėlapis:
\file{docs/markdown/en/deployment.md} ir \file{docs/markdown/en/operations.md}. Paskutinis žodis priklauso
scenarijams ir Compose failui, ir skyrius patikrintas pagal juos.


\section{Kas veikia ir ko reikia serveriui}\label{sec:ops:runs}

Visa sistema yra vienas Docker Compose projektas, aprašytas faile \file{docker-compose.yml}:
aštuoni konteineriai viename serveryje (\cref{tab:ops:services,fig:ops:host}). \Cref{ch:arch}
paaiškina, ką paslaugos daro viena kitai; čia jos išvardytos taip, kaip operatorius jas mato
komandos \texttt{docker compose ps} išvestyje.

Kiekvieno konteinerio šakninė failų sistema yra tik skaitoma, o paleidimo iš naujo politika ---
\code{unless-stopped}. Įėjimas, išėjimas, API ir duomenų bazė veikia kaip naudotojas 1000; konsolė
ir asistento paslauga veikia kaip numatytasis jų atvaizdo naudotojas (eilutė
\texttt{user: 1000:1000} joms užkomentuota); duomenų bazės įrankis ir užduočių vykdytojo
konteineris veikia kaip root. Užduočių vykdytojo konteineriui root reikia todėl, kad jis laiko
serverio Docker lizdą --- o tai daro jį root visame serveryje, --- todėl jis neturi jokio tinklo ir
nevykdo nieko, išskyrus tą vieną komandą. Compose failas nenustato nei išteklių ribų, nei
būklės patikrų, nei paleidimo tvarkos: paslaugos paleidžiamos kartu, o įėjimas atsako 502, kol
paslauga už atitinkamo kelio dar neveikia.

\begin{table}[H]
\caption{Sistemos konteineriai, kaip juos apibrėžia \file{docker-compose.yml}}\label{tab:ops:services}
\begin{tabularx}{\textwidth}{@{}L{35mm}L{33mm}Y@{}}
\toprule
\hd{Konteineris} & \hd{Atvaizdas} & \hd{Ką daro} \\
\midrule
\code{legalops-endpoint} & \code{caddy:2.11-alpine} & Vienintelis kelias į vidų, serverio 80 prievadas. \code{/api/assistant/*} siunčia į asistento paslaugą, \code{/api/*} --- į API, \code{/dbgate/*} --- į duomenų bazės įrankį, prieš tai paklausęs API, ar kviečiantysis yra administratorius, o visa kita --- į konsolę ir viešąją svetainę. Rašo JSON prieigos žurnalą. \\
\code{legalops-exit} & \code{vuknf/caddy-l4:2.11} & Vienintelis kelias į išorę: TLS ryšius persiunčia į fiksuotą serverių vardų sąrašą (\cref{sec:ops:secrets}). \\
\code{legalops-vite} & sukurtas iš \file{vite/} & Viešoji svetainė ir konsolė: Vite kūrimo serveris, o gamybos režime --- Caddy pateikiami sukompiliuoti failai. \\
\code{legalops-django} & sukurtas iš \file{django/} & API (Django). Šiame konteineryje veikia ir užduočių vykdytojo ciklas. \\
\code{legalops-worker} & \code{docker:27-cli} & Tik viena komanda, \texttt{docker exec legalops-django python3 manage.py run\_worker}: ji paleidžia užduočių vykdytojo ciklą API konteineryje ir paleidžia jį iš naujo, kai šis baigiasi. \\
\code{legalops-assistant} & sukurtas iš \file{assistant/} & Asistento paslauga: pokalbis ir kiekvienas modelio iškvietimas. \\
\code{legalops-postgres} & sukurtas iš \file{postgres/} & PostgreSQL 17.11 su sukompiliuotu pgvector plėtiniu 0.8.6. \\
\code{legalops-dbgate} & \code{dbgate/dbgate:7.3.3-alpine} & Duomenų bazės administravimo įrankis administratoriams, pasiekiamas adresu \code{/dbgate/}. \\
\bottomrule
\end{tabularx}
\end{table}

\begin{figure}[H]
\centering
\begin{tikzpicture}[x=1cm, y=1cm]
  \node[ext]   (proxy)    at (0,1.2)     {išorinis\\tarpinis serv.\\TLS};
  \node[svc]   (endpoint) at (2.75,1.2)  {įėjimas\\80 prievadas};
  \node[svc]   (vite)     at (6.1,2.45)  {vite\\svetainė ir konsolė};
  \node[svc]   (asst)     at (6.1,1.2)   {asistentas\\modelio iškvietimai};
  \node[svc]   (django)   at (6.1,-0.15) {django\\API + vykdytojo ciklas};
  \node[svc]   (worker)   at (6.1,-1.75) {užduočių vykdytojas\\tik Docker CLI};
  \node[svc]   (exit)     at (9.75,1.2)  {išėjimas\\leidžiamų adresų\\sąrašas};
  \node[store] (pg)       at (9.75,-0.15) {postgres};
  \node[svc]   (dbgate)   at (9.75,-1.75) {dbgate};
  \node[ext]   (gw)       at (12.75,1.75) {DI tinklų sąsaja};
  \node[ext]   (src)      at (12.75,0.65) {teisės aktų\\šaltiniai};

  \draw[flow] (proxy) -- (endpoint);
  \draw[flow] (endpoint) -- (asst);
  \draw[flow] (endpoint.east) -- ++(0.55,0) |- (vite.west);
  \draw[flow] (endpoint.east) -- ++(0.55,0) |- (django.west);
  \draw[flow] (endpoint.south) -- (2.75,-2.75) -| node[lbl, pos=.25] {/dbgate/, tik administratoriams} (dbgate.south);
  \draw[flow] (asst) -- (exit);
  \draw[flow] (django) -- (pg);
  \draw[flow] (dbgate) -- (pg);
  \draw[flow] ([yshift=2.5mm]django.east) -- ++(0.45,0) |- ([yshift=-2mm]exit.west);
  \draw[flow] (exit.east) -- ++(0.45,0) |- (gw.west);
  \draw[flow] (exit.east) -- ++(0.45,0) |- (src.west);
  \draw[call, <->] (django) -- node[lbl, right=1mm] {vidinės durys} (asst);
  \draw[call] (worker) -- node[lbl, right=1mm] {docker exec} (django);

  \node[store] (d1) at (1.7,-4.9)   {\_DATA/postgres\\duomenų bazė};
  \node[store] (d2) at (5.55,-4.9)  {\_DATA/django/uploads\\įkelti failai};
  \node[data]  (envf) at (9.3,-4.9) {.env\\paslaptys};
  \begin{scope}[on background layer]
    \node[lane, fit=(endpoint)(vite)(django)(asst)(dbgate)(pg)(exit)(worker), inner ysep=9pt] (stack) {};
    \node[lane, fit=(d1)(d2)(envf), fill=codebg] (disk) {};
  \end{scope}
  \node[lanelabel, anchor=south west] at (stack.north west) {vienas Docker Compose projektas: aštuoni konteineriai};
  \node[lanelabel, anchor=south west] at (disk.north west) {serverio diske, projekto aplanke};
\end{tikzpicture}
\caption{Sistema savo serveryje. Užklausos ateina tik per įėjimą, o ryšiai į išorę išeina tik per
išėjimą. Užduočių vykdytojo konteineris paleidžia užduočių vykdytojo ciklą API konteineryje ir
neturi savo tinklo. Viskas, ką sistema įsimena, guli serverio aplankuose, prijungtuose prie
duomenų bazės ir API konteinerių.}\label{fig:ops:host}
\end{figure}

\paragraph{Tinklai.} Compose failas apibrėžia \code{isolated} --- vidinį tilto tinklą
172.28.10.0/24, kuriuo sujungti visi konteineriai, išskyrus užduočių vykdytojo, ir kuris neturi
maršruto į išorę, --- ir prisijungia prie antro tinklo, \code{external}, kurį tikisi rasti jau
sukurtą: jį sukuria atnaujinimo scenarijus komanda \texttt{docker network create}, potinklyje
172.18.0.0/24, kai serveryje tokio pavadinimo tinklo nėra; serveryje, kuriame toks tinklas jau
yra, jis lieka toks, koks yra, ir sistema prie jo prisijungia. Veikiančiame serveryje šis tinklas
buvo dar iki sistemos, ir prie jo prijungti ir kiti prieglobos konteineriai (pastebėta 2026~m.
lapkričio 16~d.). Įėjimas ir išėjimas yra abiejuose tinkluose; kūrimo režime tinkle \code{external}
yra ir konsolė bei asistento paslauga, kad galėtų įsidiegti savo paketus.

\paragraph{Būsena.} Viskas, ką sistema įsimena, yra dviejuose aplankuose, kuriuos sukuria ir
naudotojui 1000 atiduoda scenarijus \file{runUpdateThisStack.sh}: \file{_DATA/postgres} --- duomenų bazės
klasteris, ir \file{_DATA/django/uploads} --- naudotojų ir administratorių įkelti failai. Paslaptys
laikomos faile \file{.env}, kurį Compose perskaito kurdamas konteinerį. Nė vieno iš šių trijų nėra
saugykloje.

\paragraph{Ką turi suteikti serveris.} Pagal \file{docs/markdown/en/deployment.md}:
\begin{itemize}
  \item Linux su Docker Engine (išbandyta su 27.3.1) ir abi Compose sąsajos: atnaujinimo
    scenarijus kviečia \code{docker-compose}, kiti scenarijai --- \texttt{docker compose}.
    Atnaujinimo scenarijus taip pat naudoja \code{sudo} ir \code{openssl}.
  \item Užtenka keturių procesoriaus branduolių ir 8~GB atminties. Nė vienam konteineriui
    neribojama atmintis; piką sukelia izoliuoti antriniai procesai, išgaunantys tekstą iš įkeltų
    failų (kiekvienam po 1,5~GB adresų erdvės, \code{CHILD_MEMORY_MB} faile
    \file{common/sandbox.py}), ir ataskaitos generavimas PDF formatu.
  \item Apie 3~GB disko vietos atvaizdams ir dar vietos duomenims.
  \item Laisvas serverio 80 prievadas ir nenaudojami du minėti potinkliai.
  \item Išeinantis HTTPS ryšys į DI tinklų sąsają (\code{ai.knf.vu.lt}) ir į tris teisės aktų
    šaltinius (\code{publications.europa.eu}, \code{www.e-tar.lt}, \code{www.lb.lt}); kol kuriami
    atvaizdai --- ir į paketų registrus.
  \item Priekyje stovintis atvirkštinis tarpinis serveris, kuris užbaigia TLS, nes įėjimas naudoja
    paprastą HTTP. Jis turi perduoti antraštę \code{Host} tokią, kokią ją siuntė naršyklė, nes
    įėjimas su ja lygina užklausos \code{Origin}, ir į antraštę \code{X-Forwarded-For}, kurią
    įėjimas nukopijuoja į \code{X-Real-IP}, įrašyti lankytojo adresą, ir tik jį (ką siunčia
    veikiančios sistemos tarpinis serveris, aprašo \cref{ch:security}). Jis neturi buferizuoti
    \code{text/event-stream} (pokalbis savo atsakymą siunčia
    srautu), turi laukti atsakymo bent 330~s (ataskaitos eksportas yra viena ilga užklausa, o
    instrukcijos bandomasis paleidimas gali laukti iki penkiolikos minučių: anksčiau jį
    nutraukiantis tarpinis serveris instrukcijų redaktoriuje parodo klaidą, nors iškvietimas vis
    tiek užbaigiamas ir įrašomas), turi praleisti 62~MiB užklausų turinį (API priima iki
    60~MiB dydžio įkeliamą teisės aktą) ir turi palikti abi įėjimo saugumo antraštes.
\end{itemize}

\paragraph{Pajėgumas.} Gamybos režime API veikia su gunicorn: \code{GUNICORN_WORKERS} procesai
po \code{GUNICORN_THREADS} gijų (numatytieji 3 ir 8, \file{django/gunicorn.conf.py}), taigi vienu
metu aptarnaujamos 24 užklausos. Vienas užduočių vykdytojo ciklas vienu metu vykdo vieną užduotį;
užduoties viduje vienu metu vyksta iki keturių modelio iškvietimų (\cref{app:config}). Užduotys iš
eilės paimamos užklausa \texttt{SELECT … FOR UPDATE SKIP LOCKED} (\code{claim} faile
\file{jobs/runner.py}), todėl antras ciklas būtų saugus; jam reikia savo \code{WORKER_PID_FILE},
nes ciklas, radęs kito ciklo PID failą, perima jo darbą, užuot veikęs šalia.
\file{docs/markdown/en/operations.md} įvardija tikrąją ribą: DI tinklų sąsajos biudžetą, ne serverį.


\section{Naujas diegimas}\label{sec:ops:install}

Diegimą sudaro keturi žingsniai šia tvarka: du vietiniai failai, vienas atnaujinimo scenarijaus
paleidimas, pirmasis administratorius ir vienas vektorių perskaičiavimas. Pati teisės aktų bazė
nei atsisiunčiama, nei modelio skaitoma iš naujo: ji atkeliauja su duomenų bazės pradine
išklotine (angl. \en{bootstrap dump}).

\subsection{Vietiniai failai ir atnaujinimo scenarijus}

\begin{codeblock}
git clone <repository> /apps/aplikacijos/comp-assistant
cd /apps/aplikacijos/comp-assistant
cp docker-compose.sample.yml docker-compose.yml
touch .env && chmod 600 .env      # then add the line  AI_GATEWAY_KEY=<the gateway key>
./runUpdateThisStack.sh
\end{codeblock}

Saugykloje versijuojamas \file{docker-compose.sample.yml}; \file{docker-compose.yml} yra vietinė
kopija, kurią vykdo diegimas, ir git ją ignoruoja, kaip ir \file{.env}.

\begin{notebox}{Dvi Compose failo eilutės, nuo kurių priklauso diegimas}
Kopija priklauso pačiam diegimui: jos nesukuria joks scenarijus, o \texttt{git pull} jos niekada
nekeičia, todėl vėlesnis pavyzdžio pakeitimas ją pasiekia tik tada, kai kas nors jį perkelia.
Naujam diegimui svarbiausios dvi pavyzdžio eilutės: tik skaitymui skirtas
\file{postgres/bootstrap.sql.gz} prijungimas prie duomenų
bazės konteinerio aplanko \file{/docker-entrypoint-initdb.d} ir įėjimo kintamasis
\code{DBGATE_AUTH_HEADER}. Be pirmojo naujas diegimas prasideda su tuščia duomenų baze:
atnaujinimo scenarijaus \code{migrate} tada sukuria schemą, o operatorius turi įkelti instrukcijas
(\code{seed_prompts}, \cref{sec:ops:prompts}) ir katalogą (\code{seed_corpus},
\cref{sec:ops:corpus}) ir vėl sumokėti už kiekvieną teisės aktų bazės skaitymą. Be antrojo įėjimas
negali prisijungti prie duomenų bazės įrankio. Saugykloje nėra failo \file{.env} pavyzdžio:
failas prasideda nuo vienos eilutės su DI tinklų sąsajos raktu, o atnaujinimo scenarijus prie jo
prideda savo sugeneruotas paslaptis.
\end{notebox}

\file{runUpdateThisStack.sh} kiekvieną kartą, taip pat ir pirmąjį, atlieka tuos pačius šešis
veiksmus:

\begin{enumerate}
  \item Sukuria \file{_DATA/django/uploads} ir \file{_DATA/postgres} ir atiduoda \file{_DATA}
    naudotojui 1000.
  \item Sugeneruoja \code{DJANGO_SECRET_KEY} (32 atsitiktinius baitus šešioliktaine forma, iš
    \texttt{openssl rand}) ir prideda jį prie \file{.env} --- tik tada, kai \file{.env} neturi
    eilutės, prasidedančios šiuo vardu. Esama eilutė, net ir tuščia, nekeičiama.
  \item Tokiu pat būdu sugeneruoja \code{ASSISTANT_INTERNAL_SECRET}. Tai bendra API ir asistento
    paslaugos vidinių durų paslaptis; jei reikšmė tuščia, abi pusės atmeta kiekvieną vidinį
    iškvietimą.
  \item Sugeneruoja \code{DBGATE_PASSWORD} --- paties duomenų bazės įrankio prisijungimo slaptažodį
    --- ir \code{DBGATE_AUTH_HEADER} --- tą patį prisijungimą, užkoduotą HTTP antraštei, kurią
    įėjimas prideda prie kiekvienos įrankiui perduodamos užklausos.
  \item Sukuria Docker tinklą \code{external} (kai tokio pavadinimo tinklas jau yra ---
    kiekvieną vėlesnį kartą arba nuo pat pradžių serveryje, kuriame jį sukūrė kas nors kita, ---
    tai klaida, bet nekenksminga), sustabdo sistemą ir vėl ją paleidžia iš naujai sukurtų atvaizdų:
    \texttt{docker-compose down}, paskui \texttt{docker-compose up -d -{}-build}.
  \item Palaukia penkias sekundes ir API konteineryje paleidžia
    \texttt{manage.py migrate -{}-noinput}.
\end{enumerate}

Scenarijus niekada nerašo \code{AI_GATEWAY_KEY}: fakulteto DI tinklų sąsajos (angl.
\en{gateway}) raktas yra vienintelė reikšmė, kurią pateikia žmogus. Be jo sistema pasileidžia,
bet asistento paslauga praneša, kad jos modelis nesukonfigūruotas, pokalbis atsako 503, o
kiekvienas užduoties etapas, kuriam reikia modelio, nepavyksta.

\subsection{Pradinė išklotinė}\label{sec:ops:bootstrap}

Kai duomenų bazės konteineris paleidžiamas su tuščiu \file{_DATA/postgres}, oficialus PostgreSQL
paleidimo scenarijus inicializuoja klasterį ir įkelia kiekvieną failą, kurį randa aplanke
\file{/docker-entrypoint-initdb.d}, --- čia vienas failas, \file{postgres/bootstrap.sql.gz}, apie
11~MB gzip suspausto paprasto SQL. Jau esamo klasterio jis niekada neliečia. Failą iš veikiančio
diegimo parašo \file{scripts/db_bootstrap_dump.sh} (\cref{fig:ops:bootstrap}).

\begin{figure}[H]
\centering
\begin{tikzpicture}[node distance=7mm and 6mm]
  \node[store] (live) {veikianti\\duomenų bazė};
  \node[svc, right=of live] (copy) {laikinas\\postgres\\be tinklo};
  \node[chk, right=of copy] (cut) {apkarpymas:\\palikti aštuonias\\lenteles, kitas\\išvalyti};
  \node[data, right=of cut] (file) {bootstrap\\.sql.gz};
  \node[store, below=13mm of live] (new) {tuščias\\duomenų aplankas};
  \node[svc, right=of new] (load) {paleidimo\\scenarijus\\įkelia failą};
  \node[man, right=of load] (admin) {pirmasis\\administratorius};
  \node[ai, right=of admin] (embed) {vienas vektorių\\perskaičiavimas};
  \draw[flow] (live) -- node[lbl, above] {skaito} (copy);
  \draw[flow] (copy) -- (cut);
  \draw[flow] (cut) -- (file);
  \draw[flow] (new) -- (load);
  \draw[flow] (load) -- (admin);
  \draw[flow] (admin) -- (embed);
  \draw[opt] (file.south) |- ([yshift=6mm]load.north) -- (load.north);
  \node[note, anchor=west] at ([xshift=3mm]file.east) {sukuriama\\veikiančioje\\sistemoje};
  \node[note, anchor=west] at ([xshift=3mm]embed.east) {naujame\\diegime};
\end{tikzpicture}
\actorlegend
\caption{Kaip naujas diegimas gauna savo teisės aktų bazę. Viršuje: \file{scripts/db_bootstrap_dump.sh}
nukopijuoja veikiančią duomenų bazę į laikiną konteinerį, ten ją apkarpo ir išsaugo rezultatą
išklotine; veikiančią duomenų bazę jis tik skaito. Apačioje: naujas diegimas vieną kartą įkelia tą
failą ir sukuria pirmąjį administratorių, kuris tada paleidžia vienintelį mokamą žingsnį ---
vektorius, kurių failas neturi.}\label{fig:ops:bootstrap}
\end{figure}

\paragraph{Ką failas turi.} Scenarijus laiko aštuonių lentelių leidžiamųjų sąrašą (masyvas
\code{keep} jo SQL kode): \code{django_migrations}, \code{legal_documents},
\code{legal_document_versions}, \code{legal_units}, \code{legal_chunks}, \code{obligations},
\code{prompts} ir \code{prompt_versions}. Todėl naujas diegimas prasideda su
\begin{itemize}
  \item schema, kurioje kiekviena migracija pažymėta kaip pritaikyta;
  \item teisės aktų baze: katalogo teisės aktais, kiekvieno ją turinčio akto dabartine redakcija,
    jos nuostatomis ir fragmentais bei viskuo, kas buvo perskaityta į redakcijų eilutes
    (apimties ir sankcijų skaitymai, aprašyti \cref{ch:obligations}, yra lentelės
    \code{legal_document_versions} stulpeliai);
  \item pareigų registru su vertintojo verdiktais ir peržiūros būsenomis --- 2026 m. lapkričio
    19 d. išleistame faile 68 teisės aktų 5\,960 pareigų;
  \item instrukcijomis su jų versijomis, kokios jos buvo išklotinės dieną: 2026~m. lapkričio
    20~d. faile yra 91 versija, pokalbio instrukcijos --- septynioliktoji.
\end{itemize}

\paragraph{Ko failas neturi.} Visos kitos lentelės išvalomos: paskyros, sesijos, įmonės ir viskas,
kas po jomis, ataskaitos, pokalbiai, programos dokumentai, užduotys, veiksmų žurnalas, modelio
iškvietimų įrašai, kainų podėlis. Kadangi sąrašas įvardija, kas lieka, vėliau į sistemą pridėta
lentelė išleidžiama tuščia, kol kas nors jos neįtraukia į tą sąrašą, --- pagal nutylėjimą niekas
nenuteka. Paliktose lentelėse scenarijus pašalina tai, kas priklauso žmonėms ar vienam diegimui:
įmonės vidaus reikalavimus (pareigas, turinčias įmonę), paliktų eilučių nuorodas į paskyras
(\code{created_by}, \code{reviewed_by}), naujesnėmis pakeistas teisės aktų redakcijas su jų vienetais ir
fragmentais ir kiekvieną vektorių. Vektoriai pašalinami dėl dydžio --- SQL tekstu jie sudarė tris
ketvirtadalius 156~MB failo, kaip rašo \file{docs/markdown/en/operations.md}, --- o kartu išvaloma ir
\code{embed_model}, todėl konsolė rodo, kurios redakcijos dar laukia savo vektorių. Neįtraukiami
ir teisės aktų šaltinių failai; kai redakcija paruošta, jų niekas nebeskaito. Apkarpymas vyksta
vienoje transakcijoje su visais atidėtais išoriniais raktais, todėl palikta eilutė, vis dar
rodanti į išvalytą lentelę, sužlugdo išklotinės kūrimą veikiančioje sistemoje, o ne diegimą,
kuris ją įkeltų.

Pagal \file{docs/markdown/en/deployment.md} įkėlimas trunka apie dešimt sekundžių
--- ilgiau nei atnaujinimo scenarijaus penkių sekundžių laukimas. Todėl pirmą kartą scenarijaus
\code{migrate} gali nepavykti su pranešimu \enquote{connection refused}; išklotinė jau
migruota, tad tokia nesėkmė tuomet nekenksminga.

\subsection{Pirmasis administratorius}\label{sec:ops:admin}

Niekas negali užsiregistruoti pats, o kiekvienam paskyrą kuriančiam maršrutui reikia
administratoriaus. Į tuščią diegimą patenkama valdymo komanda \code{create_admin}
(\file{users/management/commands/create_admin.py}):

\begin{codeblock}
read -rsp 'Admin password: ' ADMIN_PASSWORD; echo; export ADMIN_PASSWORD
docker compose exec -T -e ADMIN_PASSWORD legalops-django \
    python3 manage.py create_admin --email <address> --password-env ADMIN_PASSWORD
unset ADMIN_PASSWORD
\end{codeblock}

Saugyklos šaknyje esantis \file{runCreateAdmin.sh} paklausia adreso bei slaptažodžio ir įvykdo
būtent tai. Slaptažodis perduodamas aplinkos kintamuoju, kurį nurodo \code{--password-env}, o ne komandų
eilute, kurioje jį parodytų procesų sąrašas. Kitaip nei konsolėje nustatomam slaptažodžiui,
mažiausias ilgis jam netaikomas: slaptažodį pasirenka tas, kas serveryje vykdo komandą, o komanda
atmeta tik tuščią ir ilgesnį nei 72 bcrypt baitai (\code{PASSWORD_MAX_BYTES}). Esamai paskyrai
komanda nieko nedaro. Su \code{--reset} ji paskyrą vėl padaro aktyvia administratoriaus paskyra ir,
jei slaptažodis pateiktas, jį nustato ir užbaigia kiekvieną paskyros sesiją; taip atkuriama prieiga
po užsirakinimo, ir tai įrašoma į veiksmų žurnalą (\code{users.reset_password} arba, jei
slaptažodis nepateiktas, \code{users.update}) be veikėjo.

\subsection{Vienas vektorių perskaičiavimas}

Administratorius prisijungia, atveria teisės aktų bazės puslapį (\ui{Teisės aktų bazė}) ir jo
meniu pasirenka \ui{Perskaičiuoti visus vektorius}. Tai per \route{POST}{/api/corpus/reembed} į
eilę įrašo vieną užduotį, \code{corpus.reembed}. Užduotis (\code{reembed} faile
\file{corpus/ingest.py}) vektorizavimo modeliui po 64 (\code{EMBED_BATCH}) siunčia kiekvieno teisės
akto dabartinės redakcijos fragmentus, paskui kiekvienos pareigos lietuvišką ir anglišką tekstą,
paskui esamų įmonės dokumentų fragmentus --- naujame diegime jų nėra. Ji niekam nepriskiriama.
2026 m. lapkričio 13 d. teisės aktų bazėje buvo 4\,059 fragmentai ir 5\,960 pareigų;
\file{docs/markdown/en/deployment.md} užduotį --- apie 4\,100 fragmentų ir 11\,900 pareigų tekstų ---
vertina keliomis minutėmis ir mažiau nei vienu doleriu; tai jo įvertis, o ne datuotas matavimas.
Vektorizavimo iškvietimą aprašo \cref{ch:models}.

Kol užduotis neįvykdyta, trys dalykai matomai negerai, ir operatorius turėtų juos atpažinti kaip
praleisto žingsnio požymius, o ne kaip gedimus: teisės aktų bazės paieška randa tik raktažodžių
atitikmenis, ataskaita, kurios dokumentų rinkinys per didelis, kad būtų siunčiamas visas, neranda
įrodymų savo pareigoms (telpantis rinkinys perskaitomas visas, ir vektorių jam nereikia;
\cref{ch:report}), o teisės akto redakcijų lentelėje ten, kur
turėtų būti vektorizavimo modelio pavadinimas, rodomas brūkšnys.

Tą pačią užduotį reikia paleisti kaskart, kai pasikeičia \code{AI_EMBED_MODEL}. Dviejų modelių
vektorių palyginti negalima, todėl užduotis atnaujina visus tris vektorių rinkinius kartu; kintamasis
nustatomas dviejose vietose (API kopija įrašyta pačiame \file{docker-compose.yml}, asistento
paslaugos --- imama iš \file{.env}, kur numatytasis yra tas pats modelis), ir abiejose turi būti
nurodytas tas pats modelis.

\subsection{Rezultato patikra}

\begin{codeblock}
docker compose ps                       # every container Up
curl -s http://127.0.0.1/api/health     # {"status": "ok", "service": "legalops-django", ...}
\end{codeblock}

Pradinė išklotinė
tokia aktuali, kokia buvo paskutinį kartą ją sukuriant: kai saugyklos numatytosios instrukcijos ar katalogas yra
naujesni už failą, \code{seed_prompts} ir \code{seed_corpus} prideda tai, ko trūksta, o kas savaitę
vykdomas atnaujinimas teisės aktus padaro aktualius.


\section{Kūrimo ir gamybos režimai}\label{sec:ops:modes}

\file{docker-compose.yml} abu režimus laiko viename faile. Eilutės, priklausančios vienam režimui,
baigiasi komentaru \texttt{\# Dev} arba \texttt{\# Prod}; pateiktame faile aktyvios kūrimo režimo
eilutės. Norėdamas pereiti į gamybos režimą, operatorius užkomentuoja kiekvieną eilutę, pažymėtą
\texttt{\# Dev}, pašalina komentaro ženklą iš kiekvienos eilutės, pažymėtos \texttt{\# Prod}, ir
paleidžia \file{runUpdateThisStack.sh}. Kas pasikeičia, išvardyta \cref{tab:ops:modes}.

\begin{table}[H]
\caption{Ką vykdo du Compose failo režimai}\label{tab:ops:modes}
\begin{tabularx}{\textwidth}{@{}L{27mm}YY@{}}
\toprule
\hd{Paslauga} & \hd{Kūrimo (kaip pateikta)} & \hd{Gamybos} \\
\midrule
Konsolė ir viešoji svetainė & \file{vite/Dockerfile.dev}: \texttt{npm install} ir Vite kūrimo serveris su prijungtu \file{vite/app}; taip pat tinkle \code{external}. & \file{vite/Dockerfile}: puslapis ir jo scenarijai sukompiliuojami atvaizde (\texttt{npm ci}, vertimų sargas, \texttt{vite build}) ir pateikiami per Caddy; išeities kodas neprijungiamas, tik izoliuotas tinklas. \\
API & \code{DJANGO_RUNSERVER} yra \code{True}: Django \code{runserver} su prijungtu \file{django}, palaukus duomenų bazės. & Kintamasis nenustatytas: gunicorn su į atvaizdą įdiegtu kodu, suderintas faile \file{django/gunicorn.conf.py}. \\
Asistento paslauga & \file{assistant/Dockerfile.dev}: \texttt{npm install} ir \code{nodemon} su prijungtu \file{assistant/app}; taip pat tinkle \code{external}. & \file{assistant/Dockerfile}: \texttt{npm ci -{}-omit=dev} ir \texttt{node server.js}; 64~MB laikinoji failų sistema, tik izoliuotas tinklas. \\
\bottomrule
\end{tabularx}
\end{table}

Operatoriui svarbios dvi pasekmės. Kūrimo režime išsaugotas išeities failas iškart perkrauna
konsolę, asistento paslaugą ir API, bet ne užduočių vykdytojo ciklą: komanda
\texttt{docker compose restart legalops-worker} paleidžia ciklą su nauju kodu. Be to, kūrimo
režime asistento paslauga turi antrą tinklą su maršrutu į išorę, todėl jos srautas neapribotas
išėjimo leidžiamų adresų sąrašu; gamybos režimas tai nutraukia. Gamybos režime kodo pakeitimas
konteinerius pasiekia tik per atnaujinimo scenarijų. Ką kūrimo režimas kainuoja apsaugos
požiūriu, pasveria \cref{ch:security}.


\section{Atnaujinimas}\label{sec:ops:update}

\begin{codeblock}
git pull
./runUpdateThisStack.sh
docker compose exec legalops-django python3 manage.py seed_prompts
\end{codeblock}

Scenarijus tas pats kaip diegimo: jis sustabdo visą sistemą, iš naujo sukuria atvaizdus (dėl
podėlyje esančių sluoksnių vien kodo pakeitimas užtrunka kelias sekundes), viską paleidžia ir
migruoja duomenų bazę. Nuo \code{down} iki konteinerių grįžimo svetainė nepasiekiama; jei pirma
paleidžiama \texttt{docker compose build}, atvaizdų kūrimas į šį langą nepatenka. Trečioji eilutė
privaloma, kai atnaujinimas pakeitė failą aplanke
\file{prompts/defaults/}: \code{seed_prompts} savaime nepaleidžia niekas
(\cref{sec:ops:prompts}). Testų rinkinių scenarijus irgi nepaleidžia (\cref{ch:testing}).

Docker stabdomam konteineriui duoda dešimt sekundžių, nes Compose failas nenustato
\code{stop_grace_period}. Ką tai reiškia vykstančiam darbui:

\begin{itemize}
  \item \textbf{Pokalbio ėjimas.} Asistento paslauga iškart nustoja priimti ryšius, užklausą per
    anksčiau atvertą ir vis dar atvirą ryšį atmeta su 503 \code{SHUTTING_DOWN} ir leidžia
    vykstantiems ėjimams baigtis --- jos pačios kantrybė, \code{SHUTDOWN_GRACE_MS}, yra 190~s,
    --- bet Docker po dešimties sekundžių procesą nutraukia, ir tuo metu dar srautu siunčiamo
    ėjimo atsakymas prarandamas.
  \item \textbf{Užklausa API.} Po dešimties sekundžių dar vykdomos užklausos nutraukiamos.
  \item \textbf{Užduotis.} Užduočių vykdytojo ciklas yra antras procesas API konteineryje ir
    žūva kartu su juo. Jo vykdyta užduotis lieka pažymėta kaip vykdoma, kol kitas ciklas
    neranda, kad jos gyvybės signalas senesnis nei \code{STALE_S} (120~s, \file{jobs/runner.py});
    tada ji vėl įrašoma į eilę, jei jos rūšis leidžia dar vieną bandymą, o jei ne --- pažymima
    kaip nepavykusi su pranešimu \enquote{The worker running this job stopped unexpectedly}.
\end{itemize}

\Cref{tab:ops:jobs} rodo, kurios rūšys gauna antrą bandymą. Dauguma mokamų rūšių --- sąmoningai
--- jo negauna: pinigus išleidusi užduotis tyliai nevykdoma dukart. Todėl prieš atnaujinimą
operatorius pažiūri į konsolės užduočių puslapį (\cref{ch:console}) ir neatnaujina, kol vykdoma
ataskaita ar teisės aktų bazės skaitymas. Ataskaita, kurios užduotis taip baigiasi, pažymima kaip
nepavykusi per rūšies \code{on_abandon} kabliuką; ką išsaugo nepavykusi ataskaita, aprašyta
\cref{ch:report}.

\begin{table}[H]
\caption{Užduočių rūšys, kiek bandymų kiekvienai registruota (\code{max_attempts}) ir kas jas
paleidžia}\label{tab:ops:jobs}
\begin{tabularx}{\textwidth}{@{}L{48mm}C{16mm}Y@{}}
\toprule
\hd{Rūšis} & \hd{Bandymai} & \hd{Kas paleidžia} \\
\midrule
\code{report.analyze} & 1 & Naudotojas, pradedantis atitikties ataskaitą, ir dar kartą, kai patvirtinama jos apimtis. \\
\code{programme.generate} & 1 & Naudotojas, užsakantis programos dokumentą. \\
\code{documents.index} & 3 & Kiekvienas dokumento, turinčio tekstą, įkėlimas. \\
\code{corpus.ingest} & 2 & \code{seed_corpus}, atnaujinimas, teisės aktų bazės puslapyje pridėtas teisės aktas, įkelta akto redakcija. \\
\code{corpus.refresh_all} & 1 & Savaitinis tvarkaraštis; teisės aktų bazės puslapis. \\
\code{corpus.reembed} & 1 & Teisės aktų bazės puslapis. \\
\code{corpus.extract_obligations} & 1 & To paties pavadinimo komanda; teisės akto puslapis. \\
\code{corpus.verify_obligations} & 1 & To paties pavadinimo komanda; teisės akto puslapis. \\
\code{corpus.read_sanctions} & 1 & To paties pavadinimo komanda. \\
\code{ops.maintenance} & 1 & Dienos tvarkaraštis. \\
\bottomrule
\end{tabularx}
\end{table}

Pats savaime iš naujo paleistas užduočių vykdytojas elgiasi geriau nei nužudytas kartu su savo
konteineriu. Naujas ciklas perskaito dar veikiančio ciklo PID failą ir nusiunčia jam
\code{SIGTERM} (\code{take_over} faile \file{jobs/management/commands/run_worker.py}); senasis
ciklas nustoja imti užduotis, jo apdorojimo funkcija sustoja ties artimiausia atšaukimo patikra,
o jo užduotis grįžta į eilę \emph{neišnaudojusi} bandymo.

Baziniai atvaizdai prisegti žyma, tik ne visi vienodai tiksliai: \code{caddy:2.11-alpine},
\code{vuknf/caddy-l4:2.11} ir \code{python:3.13-slim} seka smulkiosios versijos liniją,
\code{docker:27-cli} --- pagrindinės versijos liniją, o \code{node:24.16.0-alpine} (dviejuose
kūrimo režimo atvaizduose --- \code{node:24.16.0}), \code{postgres:17.11-alpine3.24} ir
\code{dbgate/dbgate:7.3.3-alpine} nurodo tikslią laidą, todėl naujesnė jų laida atkeliauja
tik su pakeista žyma.
\file{docs/markdown/en/operations.md} prašo kartą per mėnesį prieš atnaujinimo scenarijų juos perkurti
komanda \texttt{docker compose build -{}-pull}, kad būtų įtraukta tai, kas su čia kuriamų atvaizdų
žymomis paskelbta naujai; atvaizdas, kurį Compose tik atsisiunčia, --- įėjimo, išėjimo, užduočių
vykdytojo, duomenų bazės įrankio --- atnaujinamas komanda \texttt{docker compose pull}. Nauja
PostgreSQL 17 smulkioji versija --- tai pakeista žyma faile \file{postgres/Dockerfile} ir
perkūrimas; naujai pagrindinei versijai duomenų bazę reikia išrašyti ir įkelti į tuščią naujos
versijos klasterį, ir saugykloje niekas to neautomatizuoja.


\section{Teisės aktų bazės priežiūra}\label{sec:ops:corpus}

Teisės aktų bazei reikia keturių rūšių operatoriaus dėmesio: įkelti teisės aktus, kuriais pasipildė
katalogas, palaikyti teisės aktų aktualumą, pateikti sistemai failą ten, kur ji negali jo
atsisiųsti, ir paleisti tris mokamus modelio skaitymus, kurių visai teisės aktų bazei niekas nepaleidžia
savaime (vienam aktui du iš jų administratorius gali paleisti akto puslapyje, o ataskaita pati
perskaito jai reikalingą dar neskaitytą aktą). Ką
įkėlimas ir skaitymai daro viduje, aprašo \cref{ch:corpus,ch:obligations}; šis poskyris sako, kada
juos paleisti, ką jie įrašo į eilę ir kiek kainuoja.

\subsection{Pradinis įkėlimas, atnaujinimas, failo įkėlimas}

\paragraph{Katalogo įkėlimas.} \code{seed_corpus}
(\file{corpus/management/commands/seed_corpus.py}) perskaito katalogą,
\file{corpus/catalog.yaml} --- 78 teisės aktus, iš kurių 68 įjungti, --- patikrina kiekvieną įrašą
ir sukuria dar nesamas teisės aktų eilutes. Kiekvienam naujam, įjungtam teisės aktui, kurio šaltinį
galima atsisiųsti, jis į eilę įrašo vieną \code{corpus.ingest} užduotį. Esami teisės aktai
neliečiami: po pirmojo įkėlimo tiesa yra duomenų bazė, ne failas.

\begin{codeblock}
docker compose exec legalops-django python3 manage.py seed_corpus                   # every missing act
docker compose exec legalops-django python3 manage.py seed_corpus --only mica,dora  # only these slugs
docker compose exec legalops-django python3 manage.py seed_corpus --no-ingest       # rows only
\end{codeblock}

Diegimui, pradėtam nuo pradinės išklotinės, komandos reikia tik tada, kai katalogas naujesnis už
išklotinę. Iš to, kaip ji veikia, kyla vienas perspėjimas: ji iš naujo sukuria ir teisės aktą,
kurį administratorius ištrynė sąmoningai; tai riboja \code{--only}. Įkėlimas šiek tiek išleidžia
vektorizavimui, ir tai niekam nepriskiriama: 2026 m. lapkričio 12 d. 24 naujų katalogo teisės aktų
--- 22 atsisiųstų ir dvejų gairių, kurių failą reikia įkelti (žr. toliau), --- įtraukimas
duomenų bazės kopijoje truko 89~s ir kainavo 0,08~\$.

\paragraph{Atnaujinimas.} Atnaujinimas teisės aktą iš jo šaltinio atsisiunčia iš naujo; nauja
redakcija išsaugoma tik tada, kai pasikeitė tekstas (\cref{sec:corpus:unchanged}), todėl nepakitęs
teisės aktas kainuoja
vieną užklausą ir nė vieno modelio iškvietimo. ES teisės aktas, kurio redakcijų rodyklė
neatsako, paliekamas toks, koks yra, o nesėkmė įrašoma kaip paskutinė jo klaida; nauja
redakcija, kurios tekstą teisės aktas jau yra turėjęs, perima tos ankstesnės redakcijos
pareigas ir skaitymo rezultatus (\cref{ch:corpus}). Užduočių vykdytojo tvarkaraštis kartą per savaitę į
eilę įrašo \code{corpus.refresh_all} (\cref{sec:ops:schedule}); ta užduotis įrašo po vieną
\code{corpus.ingest} kiekvienam įjungtam teisės aktui, kurio šaltinis nėra \code{upload}.
Administratorius tą patį paleidžia iš teisės aktų bazės puslapio (\ui{Tikrinti visus šaltinius};
\route{POST}{/api/corpus/refresh-all}) arba vienam teisės aktui --- iš jo puslapio
(\route{POST}{/api/corpus/documents/{id}/refresh}). Atsisiuntimo moduliai siunčia ne daugiau kaip
vieną užklausą per sekundę vienam serveriui (\code{MIN_INTERVAL_S} faile
\file{corpus/fetchers.py}) ir naudoja tik HTTPS.

\paragraph{Teisės aktas, kurio negalima atsisiųsti.} Du įjungti katalogo teisės aktai --- Europos
bankininkystės institucijos gairės dėl pinigų plovimo rizikos veiksnių ir jų pakeitimas, skirtas
kriptoturto paslaugų teikėjams, --- turi šaltinį \code{upload}: institucijos svetainės nėra
išėjimo leidžiamų adresų sąraše, todėl joms nieko neatsisiunčiama. \code{seed_corpus} sukuria jų
eilutes ir išspausdina, kad jos laukia failo. Administratorius teisės aktų bazės puslapyje atveria
kiekvieną teisės aktą, iš akto šaltinio nuorodos atsisiunčia lietuvišką PDF ir įkelia jį mygtuku
\ui{Įkelti naują redakciją}; įkeltas failas tampa laukiančia redakcija, o jo įkėlimas įrašomas į
eilę. Tas pats mygtukas tinka bet kuriam teisės aktui, kurio šaltinis neveikia, ir taip į teisės
aktų bazę patenka licencijuotas standartas (\cref{ch:corpus}). Failas gali būti iki 60~MiB
(\code{MAX_UPLOAD_BYTES} faile \file{corpus/api/views.py}).

\subsection{Trys mokami skaitymai}\label{sec:ops:readings}

Trys valdymo komandos teisės aktų bazę perduoda modeliui. Kiekviena tik \emph{įrašo} užduotis į
eilę ir iškart grįžta; skaito užduočių vykdytojas, o užduočių puslapis rodo eigą ir klaidas.
\Cref{tab:ops:readings} jas palygina, o \cref{fig:ops:readings} rodo tvarką, kuria operatorius
jas naudoja.

\begin{figure}[H]
\centering
\begin{tikzpicture}[node distance=6mm and 5mm]
  \node[svc] (ingest) {įkelti teisės aktai\\\code{seed_corpus},\\atnaujinimas, failas};
  \node[ai, right=of ingest] (extract) {\code{extract_}\\\code{obligations}};
  \node[ai, right=of extract] (verify) {\code{verify_}\\\code{obligations}};
  \node[man, right=of verify] (review) {peržiūra pareigų\\puslapyje};
  \node[ai, right=of review] (sanctions) {\code{read_}\\\code{sanctions}};
  \node[svc, below=14mm of sanctions] (dump) {\code{db_bootstrap_}\\\code{dump.sh}};
  \draw[flow] (ingest) -- (extract);
  \draw[flow] (extract) -- (verify);
  \draw[flow] (verify) -- (review);
  \draw[flow] (review) -- (sanctions);
  \draw[opt] (sanctions) -- (dump);
  \node[note, below=1mm of extract, align=center] {modelis kiekvienoje\\nuostatoje ieško pareigų};
  \node[note, below=1mm of verify, align=center] {vertintojas lygina\\kiekvieną pareigą\\su jos nuostatomis};
  \node[note, below=1mm of review, align=center] {žmogus patvirtina,\\pataiso ar atmeta};
  \node[note, left=3mm of dump, align=right] {kad nauji diegimai\\paveldėtų skaitymus};
\end{tikzpicture}
\actorlegend
\caption{Mokami teisės aktų bazės skaitymai tokia tvarka, kuria juos paleidžia operatorius.
Sankcijų skaitymas nuo kitų dviejų nepriklauso ir gali būti vykdomas bet kada; čia jis paskutinis,
nes yra pigiausias ir mažiausiai skubus.}\label{fig:ops:readings}
\end{figure}

\begin{table}[H]
\caption{Trys mokami skaitymai, kuriuos paleidžia operatorius. Kiekviena užduotis į eilę įrašoma
prioritetu $-1$ ir, paleista iš komandų eilutės, niekam nepriskiriama. Kainos --- pavieniai
datuoti paleidimai veikiančioje sistemoje su \model{gpt-6.1-sol}.}\label{tab:ops:readings}
\begin{tabularx}{\textwidth}{@{}L{25mm}YYY@{}}
\toprule
 & \hd{\code{extract_obligations}} & \hd{\code{verify_obligations}} & \hd{\code{read_sanctions}} \\
\midrule
\hd{Instrukcija} & \prompt{corpus.extract_obligations}; standartui \prompt{corpus.extract_requirements} & \prompt{corpus.verify_obligations} & \prompt{corpus.sanctions} \\
\hd{Į eilę įrašo} & Po vieną užduotį teisės aktui. & Po vieną užduotį teisės aktui. & Vieną užduotį visiems teisės aktams. \\
\hd{Be argumentų} & Kiekvieną įjungtą teisės aktą, kuris \emph{atsilikęs}: neskaitytas arba skaitytas prieš parašant dabar aktyvią instrukcijos versiją. & Kiekvieną įjungtą teisės aktą su pareigomis, kurių nematė joks vertintojas. & Kiekvieną įjungtą teisės aktą, niekada neskaitytą dėl sankcijų; standartai praleidžiami. \\
\hd{Su trumpiniais} & Tuos teisės aktus, kad ir kokia jų būsena. & Tuos teisės aktus, kiek juose yra neįvertintų pareigų; su \code{--again} --- visas jų pareigas. & Tuos teisės aktus, kiek jie niekada nebuvo skaityti; su \code{--again} --- visus. \\
\hd{Tvarka} & Pirma mažiausias teisės aktas. & Pirma mažiausiai pareigų turintis. & Pagal jurisdikciją, paskui pagal trumpinį. \\
\hd{Išmatuota} & 2026-11-10: 51 teisės aktas, 461 iškvietimas, 23,61~\$, pustrečios valandos. 2026-11-13: 17 teisės aktų, 36 iškvietimai, 4,77~\$. & 2026-11-13: 5\,960 pareigų, 369 iškvietimai, 9,80~\$. & 2026-11-13: 29 iškvietimai, 0,71~\$, 262 pakopos 27 teisės aktuose. 2026-11-19, dar kartą perskaičius kopijoje: 29 iškvietimai, 0,61~\$, 309 pakopos su antrąja instrukcijos versija; 29 iškvietimai, 0,63~\$, 288 pakopos su trečiąja. 2026-11-20 su ketvirtąja: 29 iškvietimai, 0,61~\$, 302 pakopos. \\
\bottomrule
\end{tabularx}
\end{table}

\begin{codebox}{Ką trys komandos patikrina prieš ką nors išleidžiant}
Visos trys parašytos vienodai (\file{corpus/management/commands/}). Kiekviena suranda aktyvią savo
instrukcijos versiją ir, jei jos nėra, sustoja su pranešimu \enquote{run
seed\_prompts first}. Kiekviena išspausdina instrukcijos versiją ir modelį, kuriais bus skaitomi
teisės aktai, kad ne ta instrukcija išryškėtų dar neišleidus nė cento. Kiekviena atsisako įrašyti
į eilę antrą užduotį tam pačiam teisės aktui (o sankcijų atveju --- apskritai antrą užduotį), kol
viena laukia eilėje ar vykdoma. Ir kiekviena pagal nutylėjimą ima tik tai, ko dar trūksta, todėl
po nesėkmės --- išnaudoto DI tinklų sąsajos biudžeto, sutrikimo --- vėl paleista komanda užbaigia
darbą ir už nieką nemoka dukart.
\end{codebox}

\paragraph{Prioritetai.} Užduočių vykdytojas pirmiausia paima eilėje laukiančią aukščiausio
prioriteto užduotį, o tarp lygių --- seniausią (\code{claim} faile \file{jobs/runner.py}).
Naudotojų paleisto darbo prioritetas 0 (įkelto dokumento indeksavimo --- 5), trijų komandų $-1$
(\code{PRIORITY}), o tvarkaraščio $-5$. Todėl ataskaita, įrašyta į eilę vykstant visos teisės aktų
bazės perskaitymui, vykdoma prieš skaitant kitą teisės aktą, nors ir ne anksčiau, nei baigiamas
skaityti esamas: vienas ciklas vienu metu vykdo vieną užduotį.

\paragraph{Išskyrimas --- kada teisės aktus reikia perskaityti iš naujo.} Komanda vadinasi
\code{extract_obligations}. Registro pareigos yra vienos instrukcijos versijos ir vieno modelio
darbo rezultatas. Pasikeitus išskyrimo instrukcijai ar jos modeliui, kiekvienas teisės aktas tampa
\emph{atsilikęs} (\code{is_behind} faile \file{corpus/obligations.py} palygina laiką, kada buvo
perskaityta teisės akto redakcija, su laiku, kada buvo parašyta aktyvi instrukcijos versija), ir
komanda į eilę įrašo juos visus. Atsilikę yra ir nauji teisės aktų bazės aktai, nors pirmoji
ataskaita, kuriai toks teisės aktas reikalingas, jį perskaito pati; komanda juos perskaito
anksčiau už bet kurią ataskaitą.

\begin{codeblock}
docker compose exec legalops-django python3 manage.py extract_obligations             # every act that is behind
docker compose exec legalops-django python3 manage.py extract_obligations mica dora   # these acts, whatever their state
\end{codeblock}

Masinio paleidimo užduotis turi požymį \code{only_behind}: jei, atėjus užduoties eilei, teisės aktą
jau perskaitė ataskaita, užduotis baigiasi kaip praleista, ir už nieką nemokama dukart. Teisės
aktas išlaiko esamas pasiūlytas pareigas, kol naujas jo išskyrimas sėkmingai neužbaigiamas visas, o
pareigų, kurias administratorius patvirtino ar atmetė, niekada neliečiama.

\begin{humanbox}{Peržiūra prieš perskaitymą}
Veiksmų registras (\cref{ch:register}) pareigą atpažįsta pagal jos identifikatorių, o pareiga
išlaiko savo identifikatorių, kai jos teisės aktas perskaitomas iš naujo, tik jei ją patvirtino
žmogus; pasiūlyta pareiga pakeičiama. Todėl prieš iš naujo išskirdamas teisės aktą, kurio pareigos
yra po atvirais veiksmais, operatorius pasirūpina, kad registras būtų peržiūrėtas, --- kitaip tuos
veiksmus bus galima uždaryti tik rankomis. Tai sprendimas dėl tvarkos, kurį sistema palieka savo
operatoriui.
\end{humanbox}

\paragraph{Vertintojas --- prieš registro peržiūrą.} Komanda vadinasi \code{verify_obligations}.
Vertintojas, antrasis modelio iškvietimas, kiekvieną pareigą palygina su visu jos cituojamų
nuostatų tekstu ir pažymi ją kaip pagrįstą, iš dalies pagrįstą arba nepagrįstą, pridėdamas pastabą
(\cref{ch:obligations}). Jis nieko nepatvirtina ir neatmeta: jo verdiktai tik iškelia abejotinus
įrašus pareigų puslapyje į priekį, kad žmogaus peržiūra prasidėtų ten, kur jos reikia. Komanda
paleidžiama po išskyrimo ir prieš tą peržiūrą, o su \code{--again} --- dar kartą, kai pasikeičia
vertintojo instrukcija ar modelis.

\begin{codeblock}
docker compose exec legalops-django python3 manage.py verify_obligations             # every act with unjudged obligations
docker compose exec legalops-django python3 manage.py verify_obligations --again mica  # judge these anew
\end{codeblock}

\paragraph{Sankcijos --- kiek gali kainuoti pažeidimas.} Komanda vadinasi \code{read_sanctions}.
Ataskaita pati perskaito vertinamų teisės aktų sankcijų nuostatas. Komanda skirta likusiai teisės
aktų bazei: nacionalinį įstatymą, nustatantį baudas už ES reglamento pažeidimus, išvada pagal tą
reglamentą randa tik tada, kai nacionalinis įstatymas jau perskaitytas. Ji paleidžiama vieną kartą
po to, kai pridedama teisės aktų, ir su \code{--again} --- pasikeitus instrukcijai. Vienintelė
užduotis teisės aktus skaito vieną po kito ir sustoja ties pirmuoju, kurio negalima perskaityti;
iki tol perskaityti teisės aktai išlaiko savo rezultatą, o kitas paleidimas tęsia likusius.

\begin{codeblock}
docker compose exec legalops-django python3 manage.py read_sanctions           # every enabled act never read
docker compose exec legalops-django python3 manage.py read_sanctions --again   # read all anew
\end{codeblock}

\begin{lessonbox}{Išmatuota: kiek kainuoja perskaityti visą teisės aktų bazę}
2026 m. lapkričio 10 d. visi 51 tuo metu įjungti teisės aktai buvo iš naujo perskaityti ieškant
pareigų su \model{gpt-6.1-sol}: 461 modelio iškvietimas, 23,61~\$, nuo 12:42 iki 15:12~UTC, nė
vienos nesėkmės, 4\,892 pareigos. Lapkričio 13 d. tą savaitę pridėtų 17 teisės aktų skaitymas kainavo
36 iškvietimus ir 4,77~\$ ir davė dar 1\,068 pareigas. Paskui vertintojas perskaitė visas
5\,960 pareigų per 369 iškvietimus už 9,80~\$, o visos teisės aktų bazės sankcijų nuostatoms
prireikė 29 iškvietimų ir 0,71~\$. Kartu trys 68 teisės aktų bazės skaitymai tomis kainomis
kainavo apie 39~\$ (iš jų 28,38~\$ --- išskyrimas per 497 iškvietimus) --- daugiau, nei turėjo DI
tinklų sąsajos komandos biudžetas tą dieną, kai jis
baigėsi (\cref{sec:ops:budgets}).
\end{lessonbox}

\subsection{Pradinės išklotinės atnaujinimas}

Po teisės aktų bazės atnaujinimo, išskyrimo ar peržiūros, kurią turėtų paveldėti nauji diegimai,
operatorius išklotinę sukuria iš naujo:

\begin{codeblock}
scripts/db_bootstrap_dump.sh
\end{codeblock}

Scenarijus daro tai, ką aprašo \cref{sec:ops:bootstrap}; pagal \file{docs/markdown/en/operations.md}, kai
failas buvo 9~MB, jis truko apie
40~s. Veikianti duomenų bazė tik skaitoma. Naujas failas įtraukiamas į saugyklą kartu su kodu ---
taip vieną kartą apmokėti skaitymai pasiekia kiekvieną vėlesnį diegimą.


\section{Instrukcijų priežiūra}\label{sec:ops:prompts}

Instrukcijos yra duomenų bazės eilutės, redaguojamos konsolėje (\cref{ch:prompts}); saugykla
pateikia jų numatytąsias versijas kaip YAML failus aplanke \file{prompts/defaults/} --- jų 23. Dvi
komandos sieja failus, duomenų bazę ir modelius.

\subsection{Pateiktų instrukcijų įkėlimas po atnaujinimo}

\code{seed_prompts} kviečia \code{seed_defaults} iš \file{prompts/store.py}; ši perskaito kiekvieną
numatytąjį failą, ir ją saugu paleisti kiek norima kartų:

\begin{itemize}
  \item nesama instrukcija sukuriama, o numatytoji versija tampa jos aktyvia versija;
  \item instrukcijos pavadinimas, aprašymas, kintamieji ir išvesties schema visada imami iš failo, nes
    juos skaito kodas;
  \item anksčiau nematyta numatytoji versija (atpažįstama pagal jos teksto, modelio, temperatūros
    ir išvesties ribos maišą, \code{default_hash}) tampa nauja instrukcijos versija. Ji iškart
    aktyvuojama, jei instrukcija vis dar vykdo pateiktą numatytąją versiją. Kur administratorius
    aktyvavo savo tekstą, nauja numatytoji versija laukia šalia kaip versija, kurią galima
    palyginti ir aktyvuoti rankomis;
  \item kur administratorius paliko numatytosios versijos tekstą, bet pasirinko kitą modelį ar
    kitus parametrus, naujas tekstas aktyvuojamas su tais parametrais.
\end{itemize}

Komanda išspausdina, kurios instrukcijos buvo sukurtos, kurios atnaujintos ir kurios laukia. Ji
paleidžiama rankomis po atnaujinimo, ir niekas kitas jos nepaleidžia. Pamiršus ją po
atnaujinimo, esamos instrukcijos be jokios klaidos lieka su ankstesniais tekstais. Atnaujinimo
pridėtos instrukcijos nėra, kol komanda nepaleista, kaip nėra ir jokios instrukcijos duomenų bazėje,
kuri nebuvo pradėta nuo pradinės išklotinės: užduoties etapas, kuriam jos reikia, nepavyksta su
pranešimu \enquote{Prompt '…' has no active version}.

Prieš paleidžiant komandą didelei teisės aktų bazei verta žinoti vieną šalutinį poveikį: nauja
aktyvi išskyrimo instrukcijos versija kiekvieną teisės aktą padaro \emph{atsilikusiu}, todėl kitas
\code{extract_obligations} be argumentų pasiūlo vėl perskaityti visą teisės aktų bazę
(\cref{sec:ops:readings}).

\subsection{Matavimas prieš keičiant instrukciją ar modelį}

Testų rinkiniai (\cref{ch:testing}) veikia su modelio imitacija ir įrodo, kad veikia sujungimai. Ką
instrukcijos ar jos modelio pakeitimas daro modelio atsakymams, matuoja \code{eval_prompts}
(\file{prompts/management/commands/eval_prompts.py}): ji pažymėtus atvejus pateikia tikram
modeliui per tas pačias funkcijas, kurias kviečia produktas (žymas yra uždėjęs DI modelis, žmogus
jų dar nepatvirtino). Rinkinius ir jų atvejus aprašo \cref{ch:control}; šis poskyris --- operatoriaus
pusė.

\begin{codeblock}
docker compose exec legalops-django python3 manage.py eval_prompts            # every suite
docker compose exec legalops-django python3 manage.py eval_prompts verify     # one suite
docker compose exec legalops-django python3 manage.py eval_prompts --record   # accept the scores as the baseline
\end{codeblock}

\begin{table}[H]
\caption{Šeši \code{eval_prompts} vertinimo rinkiniai ir baziniai rezultatai, įrašyti į
\file{prompts/evals/baseline.json} 2026 m. lapkričio 12 d. (sankcijų rinkinys dar kartą ---
lapkričio 20 d., ketvirtajai jo instrukcijos versijai ir su šešiolika atvejų), visi su
\model{gpt-6.1-sol}}\label{tab:ops:suites}
\begin{tabularx}{\textwidth}{@{}L{22mm}L{47mm}YL{22mm}@{}}
\toprule
\hd{Rinkinys} & \hd{Kurią instrukciją matuoja} & \hd{Ko klausia atvejis} & \hd{Bazinis rezultatas} \\
\midrule
\code{verify} & \prompt{corpus.verify_obligations} & Ar ši nuostata nustato šią pareigą? & 14 iš 14 \\
\code{scope} & \prompt{corpus.scope} & Kam taikomas paties teisės akto tekstas? & 9 iš 9 \\
\code{citations} & \prompt{corpus.check_citations} & Ar šis straipsnis reguliuoja tai, ką jam priskiria dokumentas? & 19 iš 20 \\
\code{gap} & \prompt{report.gap_analysis} & Ar pavyzdinės įmonės dokumentai atitinka šią pareigą? & 67 iš 71, leistinas nuokrypis 2 \\
\code{applicability} & \prompt{report.applicability} & Ar šis teisės aktas taikomas šiai įmonei? & 19 iš 19 \\
\code{sanctions} & \prompt{corpus.sanctions} & Kiek gali kainuoti šio straipsnio pažeidimas? & 16 iš 16 \\
\bottomrule
\end{tabularx}
\end{table}

Kiekvienam rinkiniui komanda išspausdina kiekvieną atvejį, kurį modelis atsakė neteisingai, su
tuo, ko buvo tikimasi, ir tuo, ką jis atsakė, paskui rezultatą šalia bazinio, o pabaigoje --- kiek
kainavo nuo paleidimo pradžios atlikti modelio iškvietimai (įrašytų iškvietimų suma, todėl
įskaičiuojami ir tuo pat metu vykdomo kito darbo iškvietimai). Ji baigiasi klaida, kai rinkinio
rezultatas žemesnis už bazinį. Rinkinys, kurio atsakymai nuo paleidimo iki paleidimo šiek tiek
svyruoja, savo bazinio rezultato įraše turi \code{slack} --- leistiną nuokrypį, t.~y. kiek atvejų
paleidimas gali prarasti, kol nepavyksta: spragų rinkiniui 2. Su \code{--record} rezultatai
įrašomi kaip naujas bazinis rezultatas kartu su instrukcijos versija, modeliu ir data, kurie juos
pasiekė, ir komanda nesibaigia klaida. Versija įvardijama numeriu, kurį ji turi bazinį rezultatą
įrašiusio diegimo duomenų bazėje, ir tai padaryta sąmoningai: ta pati formuluotė kitame diegime
gali turėti kitą numerį (\cref{ch:control}).

Visas šešių rinkinių paleidimas, išmatuotas 2026~m. lapkričio 12~d., kainavo apie dolerį, o beveik
pusę tos sumos, 0,42~\$, sudarė spragų rinkinys (\file{docs/markdown/en/operations.md}).
Rinkiniai \code{scope}, \code{citations}, \code{applicability} ir
\code{sanctions} skaito šio diegimo teisės aktų bazę, todėl jiems reikia teisės aktų, kuriuos
nurodo jų atvejai; atvejis, kurio teisės akto trūksta, laikomas nepavykusiu ir nurodo priežastį.
Spragų rinkinys atsineša savo įmonę: \demo{} įmonės 26 politikos ir profilis laikomi šalia atvejų aplanke
\file{prompts/evals/data/1009-ltd}, todėl kiekviename diegime jis vertina tuos pačius dokumentus.

\begin{humanbox}{Bazinio rezultato įrašymas yra žmogaus sprendimas}
Tvarka tokia: paleisti rinkinius, pakeisti instrukciją ar modelį, paleisti juos dar kartą,
perskaityti kiekvieną klaidą. Tik tada žmogus sprendžia, ar nauji rezultatai yra tas lygis,
kurio turės laikytis būsimi pakeitimai, ir juos įrašo. Klaida gali reikšti ir tai, kad neteisinga
žyma; žyma panaikinama redaguojant jos atvejo failą, o tai yra saugyklos pakeitimas, matomas jos
istorijoje.
\end{humanbox}


\section{Biudžetai ir kvotos}\label{sec:ops:budgets}

Tarp naudotojo ir DI tinklų sąsajos sąskaitos ar serverio disko stovi keturios ribos; jos
išvardytos nuo išorinės.
Kaip iškvietimas įkainojamas ir įrašomas, paaiškina \cref{ch:models}; čia --- ribos, kurias
nustato operatorius, ir kiekvienos požymiai.

\begin{table}[H]
\caption{Išlaidų ir saugojimo ribos ir kas nutinka pasiekus kiekvieną}\label{tab:ops:limits}
\begin{tabularx}{\textwidth}{@{}L{30mm}L{36mm}Y@{}}
\toprule
\hd{Riba} & \hd{Kur nustatoma} & \hd{Kas nutinka ją pasiekus} \\
\midrule
DI tinklų sąsajos komandos biudžetas & DI tinklų sąsajoje, raktui; niekas šioje sistemoje jo nekeičia. & DI tinklų sąsaja į kiekvieną iškvietimą atsako HTTP 402. Pokalbis rodo \code{AI_BUDGET_EXHAUSTED}; užduočių etapai nepavyksta su tuo pačiu kodu ir nekartojami. \\
Naudotojo dienos biudžetas & \code{LLM_DAILY_BUDGET_USD} (numatytasis 10.0; 0 reiškia neribotą) arba kiekvienai paskyrai naudotojų puslapyje (\ui{Naudotojai}). & Pokalbio ėjimas atmetamas su 429 \code{BUDGET_EXCEEDED}; taip pat instrukcijos bandomasis paleidimas, o ataskaitos ar programos dokumento pradžia --- kai lieka mažiau nei 0,25~\$. \\
Naudotojo įkėlimų kvota & \code{UPLOAD_QUOTA_FILES} (1000) ir \code{UPLOAD_QUOTA_BYTES} (2~GiB). & Įkėlimas atmetamas su 413 \code{QUOTA_EXCEEDED}. Ištrintas dokumentas atlaisvina savo dalį. \\
Vieno failo dydis & \code{MAX_FILE_BYTES}, 50~MiB, dokumentui ar pokalbio priedui; \code{MAX_UPLOAD_BYTES}, 60~MiB, teisės aktui. & Įkėlimas atmetamas. Konsolė didesnio failo atsisako dar prieš siųsdama. \\
\bottomrule
\end{tabularx}
\end{table}

\paragraph{Komandos biudžetas} yra griežta viso diegimo riba, ir ne sistemai ją didinti.

\begin{lessonbox}{Pastebėta 2026 m. rugsėjo 27 d.}
DI tinklų sąsaja atmetė kiekvieną iškvietimą su pranešimu \enquote{Team budget exceeded: 10.18 >=
10.00 dollars} (\file{docs/markdown/en/operations.md}). Tada rakto biudžetas buvo dešimt dolerių. Produktas
šią būklę visur vadina vienu kodu, \code{AI_BUDGET_EXHAUSTED}; užduoties etapas, su ja susidūręs,
iškart nepavyksta, užuot kartojęs, o teisės aktų bazės komandos pagal nutylėjimą ima tik tai, ko
dar trūksta, todėl jos nutrauktas paleidimas užbaigiamas komandą paleidus dar kartą.
\end{lessonbox}

\paragraph{Dienos biudžetas} neleidžia vienai paskyrai ištuštinti bendro rakto. Jis skaičiuojamas
kiekvienai paskyrai ir kiekvienai UTC dienai pagal įrašytą tos paskyros modelio iškvietimų kainą
(\code{budget_status} faile \file{assistant/usage.py}). Paskyros nuosava suma, jei ją nustatė
administratorius, pakeičia diegimo numatytąją; administratoriai neriboti, nebent jiems nustatyta
suma. Skaičiuojamas kiekvienas paskyrai įrašytas iškvietimas: pokalbio ėjimai, jos ataskaitų ir
programos dokumentų etapai, instrukcijų bandomieji paleidimai, modelio dalis nuorodų patikroje,
jos įkeltų failų vektorizavimas ir skaitymai, kuriuos administratorius paleidžia iš teisės akto
puslapio. Bendras darbas su teisės aktų baze niekam nepriskiriamas: pareigų išskyrimas ataskaitos
viduje, operatoriaus iš komandų eilutės paleisti skaitymai, įkėlimas ir vektorių perskaičiavimas.

Trys elgsenos pasiekus ribą yra sąmoningos. Ataskaita atmetama tik jos pradžioje ir tada, kai
patvirtinama jos apimtis; jau prasidėjęs vykdymas niekada nesustabdomas pusiaukelėje, todėl
paskyra dieną gali baigti viršijusi biudžetą vieno vykdymo kaina. Nuorodų patikra paskyrai, kuriai
nieko neliko, nesibaigia nesėkme: ji vertina antraštes pagal žodžių sutapimą, užuot klaususi modelio.
Ir įkelto failo indeksavimas taip pat nesibaigia nesėkme: užduotis \code{documents.index} pati save vėl
įrašo į eilę 00:05~UTC, praėjus penkioms minutėms po biudžeto nunulinimo (\code{RESUME_AFTER_RESET} faile
\file{companies/indexing.py}), o jos rezultatas užduočių puslapyje sako, kad ji laukia dienos
biudžeto.

\paragraph{Kvotos ir biudžetas kartu.} \file{docs/markdown/en/operations.md} apskaičiuoja blogiausią
atvejį, kurį leidžia numatytosios reikšmės: viena paskyra, užpildžiusi savo 1000 failų ir 2~GiB
kvotą tekstu, vektorizavimo modelio kainomis sukelia maždaug nuo 65~\$ iki 195~\$ vektorizavimo
išlaidų --- bet per vieną dieną niekada ne daugiau nei savo dienos biudžetą, nes indeksavimas
laukia nunulinimo. Paskyra, kuri cikle įkelia ir trina, kasdien išleidžia savo dienos biudžetą ir
ne daugiau. Tai rodo konsolės DI sąnaudų puslapis; priemonės --- mažesnis paskyros biudžetas,
mažesnė kvota arba paskyros išjungimas.

\paragraph{Prisijungimo bandymai} taip pat riboti, nors nieko nekainuoja: dešimt vienai paskyrai
per dešimt minučių (\code{EMAIL_LIMIT}, \code{EMAIL_WINDOW_S} faile \file{users/api/auth_views.py})
ir \code{LOGIN_RATE_PER_IP}, numatytasis 60, vienam kliento adresui per minutę. Veikiančiame
serveryje dauguma lankytojų ateina vienu adresu ir dalijasi tuo antruoju langu
(\cref{ch:security}). Pokalbis leidžia
\code{RATE_LIMIT_TURNS}, 12, ėjimų vienam naudotojui per minutę.


\section{Suplanuoti darbai}\label{sec:ops:schedule}

Cron konteinerio nėra. Pats užduočių vykdytojo ciklas tarp užduočių ir ne dažniau kaip kas 60~s
pažiūri į nedidelę lentelę (\code{SCHEDULE} faile \file{jobs/schedule.py},
\code{SCHEDULE_EVERY_S} faile \file{jobs/management/commands/run_worker.py}) ir įrašo rūšį į eilę,
kai naujausia tos rūšies užduotis --- kad ir kokia jos būsena ir kas ją paleido --- sukurta prieš
ilgesnį laiką nei rūšies intervalas. Rūšis, kuri dar laukia eilėje ar vykdoma, niekada neįrašoma
dukart. Suplanuotos dvi rūšys, abi prioritetu $-5$, žemiau už visa kita:

\begin{itemize}
  \item \code{corpus.refresh_all}, kas septynias dienas (\cref{sec:ops:corpus}). Administratoriaus
    rankomis paleista visų šaltinių patikra septynias dienas pradeda skaičiuoti iš naujo.
  \item \code{ops.maintenance}, kasdien (\cref{tab:ops:maintenance}).
\end{itemize}

\begin{table}[H]
\caption{Ką šalina kasdienė priežiūros užduotis (\file{ops/maintenance.py}). Užduoties
rezultatas užduočių puslapyje nurodo kiekvieno skaičių.}\label{tab:ops:maintenance}
\begin{tabularx}{\textwidth}{@{}L{38mm}YL{45mm}@{}}
\toprule
\hd{Kas} & \hd{Taisyklė} & \hd{Konstanta} \\
\midrule
Pasibaigusios sesijos & Sesijos, kurių galiojimas baigėsi. & --- (sesija gyvuoja \code{SESSION_DAYS}, 30) \\
Senos užduotys ir jų etapai & Baigtos, nepavykusios ar atšauktos daugiau nei prieš 90 dienų --- išskyrus užduotį, į kurią rodo ataskaita: ji laikoma tol, kol yra ataskaita. Dviem dalimis rengtos ataskaitos atveju tai antrosios dalies užduotis; pirmosios dalies užduotis šalinama kaip bet kuri kita (\cref{ch:arch}). & \code{JOB_RETENTION_DAYS} \\
Seni modelio iškvietimų įrašai & Senesni nei 400 dienų. Pokalbio iškvietimo įraše yra tai, kas buvo pasakyta, todėl tai kartu ir šios pokalbio turinio kopijos saugojimo terminas. & \code{TELEMETRY_RETENTION_DAYS} \\
Niekam nepriklausantys failai &Dokumentų aplanko failai, kurių nenurodo jokia dokumento eilutė ir kurie senesni nei šešios valandos: tai, ką palieka gedimas tarp failo ir eilutės įrašymo. & \code{ORPHAN_GRACE_HOURS} \\
Pasibaigę užklausų ribojimo langai & Užklausų ribų skaičiavimo langai, kurių laikas baigėsi. & --- \\
Nepanaudoti pokalbio įkėlimai & Dokumentai be įmonės, senesni nei septynios dienos, kurių nenurodo jokia pokalbio žinutė: į pokalbį įmestas, bet neišsiųstas failas. Ne daugiau kaip 500 per paleidimą. & \code{ORPHAN_UPLOAD_DAYS}, \code{ORPHAN_UPLOAD_BATCH} \\
\bottomrule
\end{tabularx}
\end{table}

Užduotis turi vieną bandymą; nepavykusį paleidimą dar kartą bando kitos dienos tvarkaraštis.


\section{Būklė, žurnalai ir gedimų šalinimas}\label{sec:ops:health}

\paragraph{Būklės patikra.} \route{GET}{/api/health} (\code{health} faile \file{ops/api/views.py})
sąmoningai vieša --- ji neneša jokių duomenų, o patikros, kuriai reikia prisijungimo duomenų,
niekas neprijungia. Ji tikrina parengtį: duomenų bazė turi atsakyti į paprastą užklausą, o į
įkeltų failų aplanką turi būti galima rašyti.

\begin{codeblock}
200  {"status": "ok", "service": "legalops-django", "database": "ok", "uploads": "ok"}
503  {"status": "error", "service": "legalops-django", "reason": "database",
      "database": "error", "uploads": "ok"}
\end{codeblock}

Atsakymas įvardija, kuri patikra nepavyko, bet niekada --- kodėl: išimties tekstas pašaliniam
atskleistų serverių vardus ir kelius, todėl priežastis yra žurnale. Asistento paslauga turi savo
patikrą, \code{/health} savo prievade izoliuotame tinkle, kuri atsako, ar sukonfigūruotas modelis,
kurie tiekėjai turi prisijungimo duomenis ir kiek pokalbio ėjimų vyksta; ji nepaskelbta per
įėjimą. Užduočių vykdytojo ciklas kiekvienos iteracijos metu ir per ilgas užduotis paliečia gyvybės
signalo failą (\file{/tmp/worker-heartbeat}); šiandien jo niekas neskaito. Aiškiausias gyvo
užduočių vykdytojo požymis --- užduotis, paliekanti eilę.

\paragraph{Žurnalai.} Kiekvienas konteineris rašo į savo standartinę išvestį;
\texttt{docker compose logs} ją skaito. Šių žurnalų nerotuoja nei Compose failas, nei
Docker tarnyba, todėl jie auga, kol konteineris pašalinamas, o tai daro kiekvienas atnaujinimas.
Įėjimas kiekvienai užklausai rašo vieną JSON eilutę su būsena ir trukme. API rašo \code{INFO} lygiu
formatu, nustatytu faile \file{settings.py}; gamybos režime gunicorn prideda prieigos eilutę. Užduočių
vykdytojo konteinerio žurnalas yra ciklo žurnalas, nes konteineris tik vykdo ciklą. PostgreSQL
registruoja sakinius, lėtesnius nei dvi sekundės (\code{log_min_duration_statement=2000}).

\begin{codeblock}
docker compose logs -f --tail 100 legalops-django     # follow the API
docker compose logs -f legalops-worker                # follow the job loop
\end{codeblock}

Viskam, kas susiję su modeliu, konsolė yra geresnis žurnalas: užduočių puslapis rodo kiekvieną
užduoties etapą su jo klaida, o DI sąnaudų puslapis --- kiekvieną įrašytą modelio iškvietimą su
jo kaina (\cref{ch:console}).

\paragraph{Dažnos nesėkmės.} \Cref{tab:ops:trouble} išvardija nesėkmes, su kuriomis operatorius
greičiausiai susidurs, --- daugumą iš \file{docs/markdown/en/operations.md}, kiekvieną patikrintą pagal kodą,
kuris sukuria pranešimą.

\begin{xltabular}{\textwidth}{@{}L{43mm}YL{42mm}@{}}
\caption{Dažnos nesėkmės, ką kiekviena reiškia ir ką daryti}\label{tab:ops:trouble}\\
\toprule
\hd{Kas matoma} & \hd{Ką tai reiškia} & \hd{Ką daryti} \\
\midrule
\endfirsthead
\toprule
\hd{Kas matoma} & \hd{Ką tai reiškia} & \hd{Ką daryti} \\
\midrule
\endhead
\bottomrule
\endlastfoot
Įėjimas atsako 502 & Paslauga už to kelio neveikia arba dar pasileidžia --- atnaujinimo metu, kol ji grįžta. & \texttt{docker compose ps}, paskui paslaugos žurnalai. \\
403 su tuščiu turiniu kiekvienai užklausai iš kitos svetainės & Įėjimas atmeta užklausą, kurios antraštė \code{Origin} nesibaigia serveriu, kuriam ji siųsta. & Taip ir turi būti. Priekyje stovintis tarpinis serveris turi perduoti naršyklės siųstą antraštę \code{Host}. \\
Atnaujinimo scenarijus baigiasi migracijos klaida & Pirmo diegimo metu duomenų bazė dar įkėlinėjo pradinę išklotinę: nekenksminga. Kitaip --- nepavyko migracija, ir sistema veikia su sena schema. & API konteineryje dar kartą paleisti \texttt{manage.py migrate}. \\
Užduoties etapas nepavyksta su \enquote{Prompt '…' has no active version} & Duomenų bazė nepradėta nuo pradinės išklotinės, ir \code{seed_prompts} niekada nepaleista. & Paleisti \code{seed_prompts}. \\
Užduotys lieka eilėje & Neveikia joks užduočių vykdytojo ciklas: užduočių vykdytojo konteineris paleidžiamas iš naujo arba neveikia API konteineris. & Užduočių vykdytojo žurnalas; \file{/tmp/worker.pid} API konteineryje. \\
Būklės patikra atsako 503 su priežastimi \code{database} & PostgreSQL neveikia, \code{DATABASE_URL} slaptažodis nėra rolės slaptažodis arba užimti visi ryšiai (\enquote{too many clients already} API žurnale). & Duomenų bazės konteinerio žurnalas; užimtus ryšius atlaisvina API konteinerio paleidimas iš naujo. \\
Būklės patikra atsako 503 su priežastimi \code{uploads}; įkėlimai nepavyksta su \enquote{Permission denied} & \file{_DATA/django} failai priklauso root. & Atiduoti aplanką naudotojui 1000; tai daro atnaujinimo scenarijus. \\
PostgreSQL vis paleidžiamas iš naujo, skųsdamasis duomenų katalogo savininku ar teisėmis & \file{_DATA/postgres} savininkas nėra naudotojas 1000 arba teisės nėra 700. & Pataisyti savininką ir teises. \\
Pokalbis atsako 503 & Nėra DI tinklų sąsajos rakto arba aktyvios \prompt{chat.assistant} versijos. & Nustatyti \code{AI_GATEWAY_KEY}; paleisti \code{seed_prompts}. \\
\code{AI_BUDGET_EXHAUSTED} pokalbyje; užduočių etapai nepavyksta su būsena 402 & Išnaudotas DI tinklų sąsajos komandos biudžetas. & Pasirūpinti, kad DI tinklų sąsajoje biudžetas būtų padidintas; tada nutrauktą darbą paleisti dar kartą. \\
\code{BUDGET_EXCEEDED} (429) & Išnaudotas paties naudotojo dienos biudžetas. & Palaukti iki 00:00~UTC arba padidinti paskyros biudžetą naudotojų puslapyje. \\
Pokalbis sako, kad asistentas paleidžiamas iš naujo (\code{SHUTTING_DOWN}) & Atnaujinimas iš naujo paleidžia asistento paslaugą. & Palaukti kelias sekundes. \\
Teisės aktas nepavyksta su \enquote{Cannot reach <host> … is it on the egress allowlist?} & Išėjimas neturi to serverio sąraše arba API konteineris neturi jam prisegto adreso. & \Cref{sec:ops:secrets}. \\
Prisijungimas atsako 429 \code{RATE_LIMITED} & Dešimt bandymų paskyrai per dešimt minučių arba pasiekta riba vienam adresui. & Palaukti tiek, kiek nurodo antraštė \code{Retry-After}. \\
Pokalbio atsakymas ateina visas iš karto arba nutrūksta & Priekyje stovintis tarpinis serveris buferizuoja srautą arba jo skaitymo laiko riba per trumpa. & \Cref{sec:ops:runs}. \\
Dokumentas lieka \ui{Teksto nėra} & Skenuotas PDF be teksto sluoksnio; sistema neturi OCR. & Įkelti PDF su tekstu arba DOCX. \\
Ataskaita, kurios dokumentų rinkinys per didelis, kad būtų siunčiamas visas, niekur neranda įrodymų; teisės aktų bazės paieška randa tik tikslius žodžius & Teisės aktų bazė niekada nebuvo vektorizuota. & Vienas vektorių perskaičiavimas (\cref{sec:ops:install}). \\
Dokumento indeksavimo užduotis sako, kad laukia dienos biudžeto & Įkėlėjo biudžetas išnaudotas; užduotis grįžta 00:05~UTC. & Nieko arba padidinti biudžetą. \\
Pakeistas \file{endpoint/Caddyfile} ar \file{exit/Caddyfile} neveikia & Konteineris vis dar vykdo seną failą. & Paleisti tą konteinerį iš naujo. \\
\end{xltabular}


\section{Pavyzdinės įmonės scenarijus}\label{sec:ops:demo}

\file{scripts/demo_1009.py} sukuria kurso pavyzdinę įmonę per viešą API, kaip naudotojas tai
darytų per konsolę, todėl kartu patikrina, ar naujas diegimas veikia nuo prisijungimo iki į eilę
įrašytos ataskaitos. Jis prisijungia su prisijungimo duomenimis, nurodytais \code{ADMIN_EMAIL} ir
\code{ADMIN_PASSWORD}, o tada:

\begin{enumerate}
  \item tarp paskyros nuosavų įmonių suranda įmonę \enquote{1009 LTD (demo)} arba sukuria ją su
    profiliu --- Lietuvos kriptoturto paslaugų teikėja, uždaroji akcinė bendrovė (UAB), kurioje
    dirba 40 darbuotojų, --- ir įrašo jos MiCA licenciją kaip leidimą su būsena \code{claimed}, su
    pastaba, kad politikos ją mini be numerio;
  \item įkelia kiekvieną pavyzdžio aplanko --- pačios saugyklos kopijos
    \file{django/compliance/prompts/evals/data/1009-ltd/documents} arba \code{MATERIALS} nurodyto
    aplanko --- Markdown failą, kurio įmonė dar neturi,
    kaip \code{policy} rūšies dokumentą (pavyzdyje jų 26: 13 apie informacijos saugumą ir 13 apie
    finansines technologijas) ir iki 15 minučių laukia, kol užduočių vykdytojas juos suindeksuos;
  \item su \code{--analyse} paprašo ataskaitos, apribotos septyniais teisės aktais, --- tais, kurių
    trumpieji pavadinimai MiCA, DORA, TFR, AMLR, PPTFPĮ, KRĮ ir GDPR, --- kainos įverčio,
    išspausdina jį ir pradeda tą ataskaitą lietuvių kalba.
\end{enumerate}

Scenarijų saugu paleisti dar kartą: jis pakartotinai naudoja įmonę ir praleidžia jau įkeltus
failus. Jam reikia tik Python su biblioteka \code{requests}, kurią turi API atvaizdas; jo antraštėje
nurodyta komanda, kuria jis paleidžiamas tame atvaizde. Numatytasis politikų aplankas yra
pačios saugyklos pavyzdžio kopija: 26 failai aplanke
\file{prompts/evals/data/1009-ltd/documents}. Indeksavimas
kainuoja centus; ataskaita kainuoja tiek, kiek išspausdina įvertis, ir abu priskiriami paskyrai, su
kuria scenarijus prisijungia (išmatuotus šios įmonės ataskaitos skaičius pateikia \cref{ch:report}).


\section{Trumpai apie paslaptis ir išėjimo leidžiamų adresų sąrašą}\label{sec:ops:secrets}

Kodėl paslaptys ir leidžiamų adresų sąrašas sutvarkyti taip, kaip sutvarkyti, paaiškina
\cref{ch:security}. Operatoriaus darbas su jais trumpas.

\paragraph{Paslapties keitimas.} Compose perskaito \file{.env} tik kurdamas konteinerį, todėl
kiekvienas keitimas baigiamas komanda \texttt{docker compose up -d}, kuri iš naujo sukuria
konteinerius, kurių reikšmės pasikeitė.

\begin{table}[H]
\caption{Paslaptys, kur kiekviena laikoma ir ką daro jos pakeitimas}\label{tab:ops:secrets}
\begin{tabularx}{\textwidth}{@{}L{42mm}L{30mm}Y@{}}
\toprule
\hd{Paslaptis} & \hd{Kur laikoma} & \hd{Keitimas} \\
\midrule
\code{AI_GATEWAY_KEY} & \file{.env} & Redaguojama rankomis, kai DI tinklų sąsaja išdavė naują raktą; asistento paslauga sukuriama iš naujo ir jį naudoja. \\
\code{ASSISTANT_INTERNAL_SECRET} & \file{.env} & Nauja atsitiktinė reikšmė; API ir asistento paslauga sukuriamos iš naujo kartu, o vidiniai iškvietimai tas kelias sekundes tarp jų nepavyksta. \\
\code{DJANGO_SECRET_KEY} & \file{.env} & Nauja atsitiktinė reikšmė. Naudotojai nieko nepastebi: sesijos yra prieigos raktai, saugomi kaip maišos, ir šiuo raktu nieko nešifruojama. \\
\code{DBGATE_PASSWORD} ir \code{DBGATE_AUTH_HEADER} & \file{.env} & Abi eilutės ištrinamos; kito paleidimo metu atnaujinimo scenarijus sugeneruoja naują porą. \\
Duomenų bazės rolės slaptažodis & \file{docker-compose.yml} & Keičiamas PostgreSQL komanda \texttt{ALTER ROLE} ir trijose vietose, kuriose jį rašo Compose failas: \code{DATABASE_URL}, \code{POSTGRES_PASSWORD} ir duomenų bazės įrankio, kuris prisijungia ta pačia role, \code{PASSWORD_con1}. \file{docs/markdown/en/operations.md} pateikia komandas, kurios naują reikšmę perduoda per aplinką, niekada ne komandų eilute. \\
Administratoriaus slaptažodis & Duomenų bazėje, maišos pavidalu & Keičiamas konsolėje; po užsirakinimo --- \texttt{create\_admin -{}-reset} (\cref{sec:ops:install}). \\
\bottomrule
\end{tabularx}
\end{table}

\paragraph{Kaip leisti jungtis prie naujo serverio.} Izoliuoto tinklo konteineris išorinį serverį pasiekia
tik tada, kai yra du dalykai: \file{exit/Caddyfile} blokas, kuris atpažįsta serverio vardą TLS
pasisveikinime ir persiunčia į to serverio 443 prievadą, ir naudojančios paslaugos
\code{extra_hosts} eilutė faile \file{docker-compose.yml}, nukreipianti vardą į išėjimo adresą
172.28.10.30. Tada \texttt{docker compose up -d} iš naujo sukuria naudojančią paslaugą, o išėjimas
paleidžiamas iš naujo, nes pasikeitė tik jo failas. Šiandien išėjimo sąraše keturi vardai: DI
tinklų sąsaja ir trys teisės aktų šaltiniai. Konsolėje užregistruotam išoriniam MCP serveriui jo
serverį reikia pridėti taip pat, kitaip jo ryšio patikra nepavyksta. Pagal
\file{docs/markdown/en/operations.md}, dviejų šaltinių sąraše nėra sąmoningai: \code{eur-lex.europa.eu} ir
\code{e-seimas.lrs.lt} į automatines užklausas atsako naršyklės patikra, kurios atsisiuntimo
modulis įveikti negali.

\paragraph{Kas paliekama operatoriui.} \file{docs/markdown/en/operations.md} baigiasi kontroliniu sąrašu
to, ko nepadengia joks kodas, --- tarp jų kliento adresas, kurį perduoda išorinis tarpinis
serveris, pačios prieglobos įrankiai, galintys skaityti projekto aplanką,
perėjimas į gamybos režimą ir duomenų bazės įrankis. \Cref{ch:security} juos aptaria kaip
likutines rizikas.
